\chapter{Definition of Foundational Theories}
\label{chap:definition-of-foundational-theories}

\section{Minimal Requirements}

To build a fundamental theory of physics, one obviously needs to know the minimum requirements
for a theory to qualify as a fundamental one.

Based on common sense, the theory must answer three basic questions:
what exists, how likely different possibilities are, and how complex structures arise from simple rules.

A candidate theory must therefore specify three ingredients: ontology, measure, and emergence.

\subsection*{1. Ontology (What exists)}

First, we must define the basic building blocks of reality.
This means specifying a well-defined space of possible states of the system.
In other words: what are the most elementary ``things'' the theory talks about?
For example, in Quantum Mechanics, the basic objects are quantum fields (or state vectors in Hilbert space) defined over spacetime.
All particles and interactions are described as excitations of these fields.

\subsection*{2. Measure (What is typical or likely)}

A theory must also explain what is typical among all possible states.
Not all states are equally relevant: some are more probable or more ``natural'' than others.
This requires a rule that assigns weights or probabilities to different states in the state space.
For example, in Quantum Field Theory, the Born rule assigns probabilities to measurement outcomes
based on the quantum state. This tells us which outcomes are likely to be observed.

\subsection*{3. Emergence (How structure arises)}

Finally, a fundamental theory must explain how complex structures arise from simple underlying rules.
This is the bridge between microscopic laws and the macroscopic hierarchical structures we observe.
For example, chemistry emerges from Quantum Field Theory.
atoms form molecules, and molecules form complex structures, even though the underlying description
is entirely in terms of quantum fields and their interactions.

\subsection*{Hierarchy}

Without ontology, the theory has nothing to describe. Without a measure, the theory cannot make predictions.
Without emergence, the theory cannot connect to the macroscopic hierarchical structures we observe.

As with any well-designed modular software, the design should also be mutually orthogonal without cyclic dependencies.
Ontology defines the domain. Measure provides weighting on that domain, so it depends on the ontology.
Emergence depends on both ontology and measure. This yields a clean layered design free of cyclic dependencies.

\textbf{The dependency stack:}

\[
\begin{array}{c}
\mathbf{[Emergence]} \\
\downarrow \ \text{\small{Depends on both}} \\
\mathbf{[Measure]} \\
\downarrow \ \text{\small{Depends on Ontology}} \\
\mathbf{[Ontology]} \\
\end{array}
\]

Of the three pillars, emergence is arguably the least fundamental.
Given a sufficiently rich ontology together with a suitable measure or dynamical law, higher-level structure may arise automatically.
In that sense, emergence can often be viewed as a consequence rather than an independent primitive.



\section{Classifying Current Theories}

Many candidate theories implement the three ingredients to varying degrees of completeness.

\subsection*{Ontology (What Exists)}
\begin{itemize}
    \item \textbf{Physical primitives:} Particles, fields, spacetime (e.g., Standard Model, General Relativity)
    \item \textbf{Information-theoretic:} Binary distinctions, bits (e.g., Wheeler's ``It from Bit'')
    \item \textbf{Computational:} Programs, Turing machines or cellular automata
    \item \textbf{Mathematical:} Abstract mathematical structures (e.g., Tegmark's Mathematical Universe Hypothesis)
\end{itemize}

\subsection*{Measure (Typicality)}
\begin{itemize}
    \item \textbf{Explicit probabilistic measure:} Classical statistical mechanics (e.g., Boltzmann distributions)
    \item \textbf{Algorithmic measure:} Program-weighted probability (e.g., Solomonoff prior in Algorithmic Information Theory)
    \item \textbf{Implicit or undefined measure:} Many multiverse or mathematical universe proposals
    \item \textbf{Emergent measure:} Measure derived from structural or dynamical properties
\end{itemize}

\subsection*{Emergence}
\begin{itemize}
    \item \textbf{Dynamical laws:} Evolution governed by differential equations (e.g., GR, QFT)
    \item \textbf{Statistical emergence:} Order arising from fluctuations and typicality (e.g., Boltzmann)
    \item \textbf{Computational emergence:} Structure arising from execution of simple programs or rules (e.g., cellular automata)
    \item \textbf{Undefined or implicit emergence:} Present in many abstract frameworks without explicit mechanism
\end{itemize}


Then there is String theory / M-theory, Loop Quantum Gravity, Quantum information, and many more 
somewhat blurring the line between ontology and measure.

Even with interpretations of Many-Worlds or holographic approaches the separation between ontology and measure can become fuzzy as the measure may appear partly emergent from the dynamics or entanglement structure.

Whoever designed the universe might have been in a hurry.


\section{Conclusions}

To map this onto our earlier dilemma about reality: ontology is our attempt to define what is ``Real''---Penrose's fundamental Platonic reality. Conversely, measure and emergence are the ``Abstract'' mathematical machinery used to calculate predictions---Hawking's pragmatic toolbox.

If we accept the pragmatic path (\textbf{maximized ontology}), we accept a complex, highly specific starting state of fields, strings, or dimensions without explaining why they exist.
The reward is incredibly precise predictions.
The cost, is that this path can not lead to a true Theory of Everything, because it relies on a mountain of unexplained ``magic'' baseline assumptions.

If we choose the ultimate path (\textbf{minimized ontology}), we gain the potential to explain absolutely everything without arbitrary starting conditions.
But the price to pay is maximized abstractness and a completely dissolved sense of physical reality. 

By minimizing ontology we maximize abstractness, which points toward a unified, complete theory. By maximizing ontology, we imply there is no ultimate Theory of Everything---only a list of highly successful, unexplained assumptions. 

Does the universe have a solid foundation, or is it abstract all the way down?
